\documentclass[11pt,a4paper]{article}
\usepackage[margin=1in]{geometry}
\usepackage{amsmath,amssymb,xcolor}
\usepackage{listings}
\usepackage{graphicx}

\setlength{\parindent}{0pt}
\newcommand{\solution}{\par\medskip
    \textcolor{blue}{\textbf{Solution:}}\par\smallskip}

\begin{document}

\begin{center}
    {\Large\bfseries SI152 --- Fall 2026 --- Homework 3}\\[0.6em]
    \textbf{Due: 2026/11/1, 23:59}
\end{center}

\textbf{Submission:} Submit a single PDF containing your written
solutions, completed functions, and required plots.
No notebook submission is required.


\section*{Problem 1 (Conjugate Gradient Method)}

Let $\mathbf{A}\in\mathbb{R}^{n\times n}$ be symmetric positive
definite. Consider solving $\mathbf{A}\mathbf{x}=\mathbf{b}$ by CG,
with $\mathbf{r}_k=\mathbf{A}\mathbf{x}_k-\mathbf{b}$.
Assume exact arithmetic and stop when $\mathbf{r}_k=\mathbf{0}$.

\begin{enumerate}
    \item Prove that any collection of nonzero, pairwise
    $\mathbf{A}$-conjugate vectors is linearly independent.
    Hence, show that such a collection contains at most $n$ vectors.
    \textbf{(5 points)}

    \item Prove that the search directions generated by preliminary
    CG (slide 42) are pairwise $\mathbf{A}$-conjugate. 
    \textit{Hint: Use induction and
    $\mathbf{A}\mathbf{p}_j
    =(\mathbf{r}_{j+1}-\mathbf{r}_j)/\alpha_j$.}
    \textbf{(15 points)}

    \item Show that preliminary CG (slide 42) and practical CG
    (slide 46), starting from the same $\mathbf{x}_0$, generate
    the same iterates.
    \textbf{(15 points)}

    \item Apply practical CG (slide 46) to
    \[
        \mathbf{A}=
        \begin{bmatrix}
            2&1&1\\
            1&2&1\\
            1&1&2
        \end{bmatrix},
        \qquad
        \mathbf{b}=\begin{bmatrix}1\\1\\0\end{bmatrix},
        \qquad
        \mathbf{x}_0=\mathbf{0}.
    \]
    Compute $\alpha_0$, $\mathbf{x}_1$, $\mathbf{r}_1$,
    $\beta_1$, $\mathbf{p}_1$, $\alpha_1$, $\mathbf{x}_2$,
    and $\mathbf{r}_2$ by hand, and verify that CG terminates
    after exactly two iterations.
    \textbf{(10 points)}
\end{enumerate}

\solution
% Write the solution here.

\begin{enumerate}
    \item[1. ]
    
    Let $\{p_1, p_2, \dots, p_k\}$ be a collection of nonzero, pairwise $\mathbf{A}$-conjugate vectors.

    Construct $\sum_{i=1}^k \alpha_i p_i = 0$ for some $\alpha_i$.

    Take $j \in \{1, 2, \dots, k\}$ and multiply by $p_j^T A$:

    $$
    p_j^T A \left( \sum_{i=1}^k \alpha_i p_i \right) = \sum_{i=1}^k \alpha_i (p_j^T A p_i) = 0
    $$

    $$
    \alpha_j (p_j^T A p_j) = 0
    $$

    $A \succ 0$, so $p_j^T A p_j > 0$, hence $\alpha_j = 0$.

    As a result, $\{p_1, p_2, \dots, p_k\}$ is linearly independent.

    And since the dimension of $\mathbb{R}^n$ is $n$, any linearly independent set can have at most $n$ vectors.

    \item[2. ]
    
    $r_1 = b - A x_1$, $p_1 = r_1$.

    $$
    \begin{aligned}
    \alpha_k &= \frac{p_k^\top r_k}{p_k^\top A p_k} \\
    x_{k+1} &= x_k + \alpha_k p_k \\
    r_{k+1} &= b - A x_{k+1} \\
    \beta_{k+1} &= \frac{p_k^\top A r_{k+1}}{p_k^\top A p_k} \\
    p_{k+1} &= r_{k+1} - \beta_{k+1} p_k
    \end{aligned}
    $$

    $$
    r_{k+1} = b - A x_{k+1} = b - A (x_k + \alpha_k p_k) = r_k - \alpha_k A p_k
    $$

    When $k = 1$, $p_1 = r_1$.

    $$
    p_2^\top A p_1
    = (r_2 - \beta_2 p_1)^\top A p_1
    = r_2^\top A p_1 - \beta_2 p_1^\top A p_1
    = r_2^\top A p_1 - \frac{p_1^\top A r_2}{p_1^\top A p_1} p_1^\top A p_1
    =0
    $$

    Assume $p_i^\top A p_j = 0$ for all $i < j \leq k$. Then for $p_{k+1}$:

    When $i = k$,

    $$
    p_{k+1}^\top A p_k
    = (r_{k+1} - \beta_{k+1} p_k)^\top A p_k
    = r_{k+1}^\top A p_k - \beta_{k+1} p_k^\top A p_k
    = p_k^\top A r_{k+1} - \frac{p_k^\top A r_{k+1}}{p_k^\top A p_k} (p_k^\top A p_k)
    = 0
    $$

    When $i < k$,

    $$
    p_{k+1}^\top A p_i
    = (r_{k+1} - \beta_{k+1} p_k)^\top A p_i
    = r_{k+1}^\top A p_i - \beta_{k+1} p_k^\top A p_i
    $$

    $p_k^\top A p_i = 0$, so $p_{k+1}^\top A p_i = r_{k+1}^\top A p_i$.

    $$
    r_{k+1}^\top A p_i = \frac{1}{\alpha_i} r_{k+1}^\top (r_i - r_{i+1}) = \frac{1}{\alpha_i} (r_{k+1}^\top r_i - r_{k+1}^\top r_{i+1})
    $$

    $$
    r_{k+1} = r_{j+1} - \sum_{m=j+1}^k \alpha_m A p_m
    $$

    $$
    p_j^\top r_{k+1} = p_j^\top r_{j+1} - \sum_{m=j+1}^k \alpha_m p_j^\top A p_m
    $$

    When $j < m \le k$, $p_j^\top A p_m = 0$. Therefore,

    $$
    p_j^\top r_{k+1} = p_j^\top r_{j+1}
    $$

    $$
    p_j^\top r_{j+1} = p_j^\top (r_j - \alpha_j A p_j) = p_j^\top r_j - \frac{p_j^\top r_j}{p_j^\top A p_j} (p_j^\top A p_j) = 0
    $$

    Hence, $p_j^\top r_{k+1} = 0$ for all $j$.

    $p_m = r_m - \beta_m p_{m-1}$, so $r_m \in \mathrm{span}\{p_1, p_2, \dots, p_m\}$.

    Therefore,

    $$
    r_{k+1}^\top r_i = 0,\quad r_{k+1}^\top r_{i+1} = 0
    $$

    As a result, $p_{k+1}^\top A p_i = 0$ for all $i < k+1$.

    According to the principle of mathematical induction, the search directions generated by preliminary CG are pairwise $\mathbf{A}$-conjugate.

    \item[3. ]
    
    Assuming preliminary CG and practical CG share the same first $k$ steps.

    $$
    \alpha_k
    = \frac{p_k^\top r_k}{p_k^\top A p_k}
    = \frac{(r_k + \beta_k p_{k-1})^\top r_k}{p_k^\top A p_k}
    = \frac{r_k^\top r_k + \beta_k p_{k-1}^\top r_k}{p_k^\top A p_k}
    = \frac{\Vert r_k \Vert ^2}{p_k^\top A p_k}
    $$

    $$
    r_k - \alpha_k A p_k = (b - A x_k) - A(\alpha_k p_k) = b - A(x_k + \alpha_k p_k) = b - A x_{k+1}
    $$

    $$
    p_k^\top Ar_{k+1} = r_{k+1}^\top Ap_k = r_{k+1}^\top \left( \frac{r_k - r_{k+1}}{\alpha_k} \right) = \frac{r_{k+1}^\top r_k - \Vert{}r_{k+1}\Vert{}^2}{\alpha_k}
    $$

    $$
    p_k^\top Ar_{k+1} = -\frac{\Vert{}r_{k+1}\Vert{}^2}{\alpha_k}
    $$

    $$
    p_k^\top Ap_k = \frac{p_k^\top r_k}{\alpha_k} = \frac{\Vert{}r_k\Vert{}^2}{\alpha_k}
    $$

    $$
    \beta_{k+1}^{\text{pre}} = \frac{-\Vert{}r_{k+1}\Vert{}^2 / \alpha_k}{\Vert{}r_k\Vert{}^2 / \alpha_k} = -\frac{\Vert{}r_{k+1}\Vert{}^2}{\Vert{}r_k\Vert{}^2} = -\beta_{k+1}
    $$

    According to the principle of mathematical induction, preliminary CG and practical CG generate the same iterates.

    \item[4. ]
    
    $$
    A
    =
    \begin{bmatrix}
        2 & 1 & 1 \\
        1 & 2 & 1 \\
        1 & 1 & 2
    \end{bmatrix}
    , \quad
    b
    =
    \begin{bmatrix}
        1 \\
        1 \\
        0
    \end{bmatrix}
    , \quad
    x_0
    =
    \begin{bmatrix}
        0 \\
        0 \\
        0
    \end{bmatrix}
    $$

    $$
    r_0 = b - A x_0
    =
    \begin{bmatrix}
        1 \\
        1 \\
        0
    \end{bmatrix}
    , \quad
    p_0 = r_0
    =
    \begin{bmatrix}
        1 \\
        1 \\
        0
    \end{bmatrix}
    $$

    $$
    \Vert r_0 \Vert ^2 = 2
    $$

    $$
    A p_0
    =
    \begin{bmatrix}
        2 & 1 & 1 \\
        1 & 2 & 1 \\
        1 & 1 & 2
    \end{bmatrix}
    \begin{bmatrix}
        1 \\
        1 \\
        0
    \end{bmatrix}
    =
    \begin{bmatrix}
        3 \\
        3 \\
        2
    \end{bmatrix}
    $$

    $$
    p_0^\top A p_0
    = 
    \begin{bmatrix}
        1 & 1 & 0
    \end{bmatrix}
    \begin{bmatrix}
        3 \\
        3 \\
        2
    \end{bmatrix}
    = 6
    $$

    $$
    \alpha_0
    = \frac{\Vert  r_0 \Vert ^2}{p_0^\top A p_0} 
    = \frac{1}{3}
    $$

    $$
    x_1
    = x_0 + \alpha_0 p_0
    = \begin{bmatrix}
        1/3 \\
        1/3 \\
        0
    \end{bmatrix}
    $$

    $$
    r_1
    = r_0 - \alpha_0 A p_0
    = \begin{bmatrix}
        0 \\
        0 \\
        -2/3
    \end{bmatrix}
    $$

    $$
    \Vert{} r_1 \Vert{}^2 = \frac{4}{9}
    $$

    $$
    \beta_1
    = \frac{\Vert r_1 \Vert ^2}{\Vert r_0 \Vert ^2}
    = \frac{2}{9}
    $$

    $$
    p_1
    = r_1 + \beta_1 p_0
    = \frac{2}{9}\begin{bmatrix} 1 \\ 1 \\ -3 \end{bmatrix}
    $$

    $$
    A p_1
    =
    \begin{bmatrix}
        0 \\
        0 \\
        - \frac{8}{9}
    \end{bmatrix}
    $$

    $$
    p_1^\top A p_1
    = \frac{16}{27}
    $$
    
    $$
    \alpha_1 = \frac{3}{4}
    $$

    $$
    x_2 = \frac{1}{2} \begin{bmatrix} 1 \\ 1 \\ -1 \end{bmatrix}
    $$

    $$
    r_2 = b - A x_2
    = \begin{bmatrix} 0 \\ 0 \\ 0 \end{bmatrix}
    $$

    A has two distinct eigenvalues, so CG terminates after at most two iterations.
\end{enumerate}

\newpage



\section*{Problem 2 (CG for Sparse Least Squares)}

Consider the regularized least-squares problem
\[
    \min_{\mathbf{x}\in\mathbb{R}^n}
    \frac12\|\mathbf{B}\mathbf{x}-\mathbf{b}\|_2^2
    +\frac{\lambda}{2}\|\mathbf{x}\|_2^2,
\]
where $\mathbf{B}\in\mathbb{R}^{m\times n}$ is sparse and
$\lambda>0$. Use the provided notebook for implementation
and experiments.

\begin{enumerate}
    \item Derive the normal equations and prove that their
    coefficient matrix is symmetric positive definite.
    \textbf{(5 points)}

    \item Complete the \texttt{TODO} sections in the notebook,
    keeping the function interfaces unchanged.
    Implement practical CG and the matrix-vector product
    without explicitly forming $\mathbf{B}^\top\mathbf{B}$.
    Run the experiments with the provided parameters.
    In your written answer, include the completed
    \texttt{matvec}, \texttt{cg}, and \texttt{run\_cg} functions,
    and the timing plot generated by the final code cell.
    \textbf{(20 points)}

    \item Explain in 1--2 sentences why the CG implementation
    accepts a \texttt{matvec} function rather than an explicit
    matrix. What computational and storage advantages does
    this provide for the problem above?
    \textbf{(5 points)}
\end{enumerate}

\solution

\begin{enumerate}
    \item[1. ]
    
    $$
    \nabla f(x) = B^\top (Bx - b) + \lambda x = 0
    $$

    $$
    (B^\top B + \lambda I) x = B^\top b
    $$

    $$
    (B^\top B + \lambda I)^\top = B^\top B + \lambda I
    $$

    $\forall x \in \mathbb{R}^n$, 

    $$
    x^\top (B^\top B + \lambda I) x
    = x^\top B^\top B x + \lambda x^\top x
    = \Vert Bx \Vert _2^2 + \lambda \Vert x \Vert_2^2 > 0
    $$

    As a result, $B^\top B + \lambda I$ is symmetric positive definite.

    \item[2. ]
    
    \begin{lstlisting}[
        language=Python,
        basicstyle=\ttfamily\footnotesize,
        breaklines=true,
        frame=single,
        showstringspaces=false
    ]
    def matvec(B, v, lam):
    """Inputs: sparse B (m,n), vector v (n,), scalar lam > 0.
    Output: vector (B.T @ B + lam*I) @ v, shape (n,).
    """
    v = (B.T @ (B @ v)) + lam * v
    return v
    \end{lstlisting}

    \begin{lstlisting}[
        language=Python,
        basicstyle=\ttfamily\footnotesize,
        breaklines=true,
        frame=single,
        showstringspaces=false
    ]
    def cg(matvec, b, epsilon=1e-8, maxiter=2000):
    """Inputs: callable matvec(v) returning A @ v, vector b (n,),
    relative tolerance epsilon, and integer iteration limit maxiter.
    Outputs: solution x (n,) and integer iteration count.
    The matrix A is symmetric positive definite.
    """
    norm_b = np.linalg.norm(b)
    if norm_b == 0:
        return np.zeros_like(b), 0

    x = np.zeros_like(b)
    r = b.copy()
    p = r.copy()
    
    r_dot_r = np.dot(r, r)
    stop_tol = (epsilon * norm_b) ** 2

    if r_dot_r < stop_tol:
        return x, 0

    for k in range(maxiter):
        Ap = matvec(p)
        alpha = r_dot_r / np.dot(p, Ap)
        x = x + alpha * p
        r = r - alpha * Ap

        r_dot_r_next = np.dot(r, r)
        if r_dot_r_next < stop_tol:
            return x, k + 1

        beta = r_dot_r_next / r_dot_r
        p = r + beta * p
        r_dot_r = r_dot_r_next

    return x, maxiter
    \end{lstlisting}

    \begin{lstlisting}[
        language=Python,
        basicstyle=\ttfamily\footnotesize,
        breaklines=true,
        frame=single,
        showstringspaces=false
    ]
    def run_cg(B, b, lam, epsilon, maxiter):
        """Return the CG solution and iteration count for this least-squares problem."""
        mv = lambda v: matvec(B, v, lam)
        rhs = B.T @ b
        return cg(mv, rhs, epsilon=epsilon, maxiter=maxiter)
    \end{lstlisting}

    \begin{figure}[h]
        \centering
        \includegraphics[width=0.8\textwidth]{timing_plot_CG.png}
        \caption{Timing plot for CG implementation.}
    \end{figure}

    \item[3. ]
    
    Every iteration of CG only requires $Av$ without knowing $A$ explicitly. Hence, a matvec function can decouple the algorithm logic and the specific implementation of operator.

    B is a sparse matrix, so storing $B^T B$ explicitly would require more memory and time. And using matvec to compute $B^T B v$ only takes $O(nnz(B))$ time.

    What's more, explicitly forming $B^T B$ may destroy sparsity and lead to enormous preprocessing computation and numerical instability, while using matvec preserves the structure of the problem.
\end{enumerate}

\newpage



\section*{Problem 3 (Levenberg--Marquardt Method)}

Consider the nonlinear least-squares problem
\[
    \min_{\mathbf{x}\in\mathbb{R}^n}
    \frac12\sum_{i=1}^{n-1}
    \left[100(x_{i+1}-x_i^2)^2+(1-x_i)^2\right].
\]
Complete the \texttt{TODO} section in the provided notebook
to implement LM using the specified Madsen--Nielsen rules.
Reuse your CG code from Problem 2 to solve each subproblem.
Run the provided experiments for $n=1000$.
In your written answer, include the completed \texttt{lm}
function and the final convergence plot.
\textbf{(25 points)}

\solution
Write the solution here.

\begin{lstlisting}[
    language=Python,
    basicstyle=\ttfamily\footnotesize,
    breaklines=true,
    frame=single,
    showstringspaces=false
]
def lm(x0, tau=1e-3, tol=1e-6, maxiter=5000, cg_tol=1e-6, cg_maxiter=2000):
    """Inputs: initial vector x0; positive damping scale tau;
    outer/inner tolerances and iteration limits.
    Outputs: x (n,) and objective history (initial value plus every attempt).
    """
    x = np.array(x0, dtype=float, copy=True)
    r, J = residual_jacobian(x)
    f = 0.5 * (r @ r)
    history = [f]

    diag_JtJ = np.asarray(J.power(2).sum(axis=0)).ravel()
    mu = tau * np.max(diag_JtJ)
    nu = 2.0

    g = J.T @ r

    for _ in range(maxiter):
        if np.linalg.norm(g, ord=np.inf) <= tol:
            break

        def matvec(v):
            return J.T @ (J @ v) + mu * v

        d, _ = cg(matvec, -g, epsilon=cg_tol, maxiter=cg_maxiter)

        Jd = J @ d
        predicted_reduction = -np.dot(g, d) - 0.5 * np.dot(Jd, Jd)

        x_trial = x + d
        r_trial, J_trial = residual_jacobian(x_trial)
        f_trial = 0.5 * np.dot(r_trial, r_trial)

        if predicted_reduction > 0 and np.isfinite(f_trial):
            rho = (f - f_trial) / predicted_reduction
        else:
            rho = -1.0

        if rho > 0:
            x = x_trial
            r = r_trial
            J = J_trial
            f = f_trial
            g = J.T @ r

            mu = mu * max(1.0 / 3.0, 1.0 - (2.0 * rho - 1.0) ** 3)
            nu = 2.0
        else:
            mu = mu * nu
            nu = 2.0 * nu

        history.append(f)

    return x, np.array(history)
\end{lstlisting}

\begin{figure}[h]
    \centering
    \includegraphics[width=0.8\textwidth]{timing_plot_LM.png}
    \caption{Timing plot for LM implementation.}
\end{figure}


\end{document}